\documentclass[../main.tex]{subfiles}
\usepackage{amssymb}
\usepackage{amsmath,amsthm} 
\usepackage{amsfonts} 
\usepackage{graphicx} 
\graphicspath{{\subfix{../pictures/}}}
\begin{document}
\section{Uživatelská příručka}
Pro spuštění programu je potřeba mít nainstalovanou sadu kompilátorů gcc a to jak pro Linux, tak i pro Windows. Na Linuxu je třeba pro zkompilování zavolat příkaz \textbf{"make"} nebo \textbf{"make \textless filename\textgreater"}, kde \textless filename\textgreater je název makefile k použití. Pro Windows je to obdobné s rozdílem že místo příkazu \textbf{"make"} se použije \textbf{"mingw32-make"}.
\vspace{0.5em}

Nyní, když je program zkompilovaný, tak je možn ho spustit. Na Windows je potřeba otevřít konzoli nebo PowerShell a dostat se do adresáře se zkompilovaným souborem. Pro zpuštění v konzoli se zadá příkaz: \newline

\textbf{C:$\backslash\textgreater$ graph.exe $\textless$ function $\textgreater \textless$ filename$\textgreater[\textless$ limits$\textgreater]$}
\vspace{1em}

Kde \textless function\textgreater \ je matematická funkce, \textless filename\textgreater je název souboru (například test.ps) a \textless limits\textgreater je nepovinný parametr limit. (ro PowerShell to je obdobné akorát je třeba na začátku přidat \textbf{"./"}.
Pro Linux je spouštění souboru stejné jako u PowerShellu, tedy:
\vspace{1em}

\textbf{C:$\backslash\textgreater$ ./graph.exe $\textless$ function $\textgreater \textless$ filename$\textgreater[\textless$ limits$\textgreater]$}
\vspace{1em}

Po zadání spouštěcího příkazu se program rozběhne a pokud nevypíše chybové hlášení a skončil s kódem 0, tak vše proběhlo v pořádku a soubor byl vytvořen. Pokud ale byla v zadaných parametrech chyba, tak program vypíše co je za problém a skončí s chybovým kódem dle zadání. Výsledky je možné vidět na stránce \pageref{ref_images}
\vspace{0.5em}

Zde je ukázka správně zapsaného Spouštěcího příkazu:
\vspace{1em}

\textbf{C:$\backslash\textgreater$ ./graph.exe  sin(xˆ2)*cos(x) test.ps 6.28:12.56:-1:1}
\vspace{1em}
\end{document}